\section{Detailed benchmark selection and SOTA identification}
\label{app:collection}

This appendix documents the protocol and the artifact locations cited by the working draft.
Paths are relative to the evaluation repository, whose release URL remains \todo{to be added}.
The quantitative audit snapshots below are dated 2026-09-11 unless stated otherwise. They
are historical evidence about the pipeline and must not be interpreted as the current completion
status of the ongoing 25\%-subset comparator evaluations. The later protocol is governed by
\texttt{experiment\_design.md} (2026-09-12) and the dated rulings in
\texttt{4\_comparators/DECISIONS.md}; these supersede conflicting earlier contract prose.

\paragraph{Collection artifacts and taxonomy.}
\texttt{1\_collection/benchmark\_audit\_full.csv} contains the candidate records and
\texttt{1\_collection/ready\_manifest.json} the archived 83-entry evidence shortlist. The filename
is an implementation label: it does not certify execution readiness. The 30-task inventory
records candidates brought into feasibility and task-fidelity review. These dated artifacts
require reconciliation of their memberships and inclusion decisions before they can support
exact exclusion counts between stages. The 554 records carry
624 assignments across 33 overlapping domain labels. The largest labels are reasoning (93),
agent (59), medical (38), math (37), multimodal (32), retrieval (28), web (26), clinical (26),
chemistry (25), vision (22), biology (22), and code (20). These are candidate-pool labels, not the
domain distribution of the updated evaluation suite. Coding-labeled candidates are included
in the audit denominator but excluded from the non-coding deployment scope.

\paragraph{Screening gates.}
The scope gate checks whether a benchmark belongs to the non-coding study. The legitimacy gate
records whether a usable comparator exists and whether independent papers have evaluated the
benchmark. A comparator reported only by the benchmark's authors is marked
\texttt{author\_reported\_only}, rather than silently discarded. The saturation gate checks
whether a reproducible comparator and useful headroom remain. T1 requires a non-saturated
benchmark with a reproducible multi-agent reference system; T2 admits more limited
reproducibility or headroom; T3 has no agent reference above the bare-model reference; SKIP
contains the remaining excluded cases. The historical evidence tiering has 122 T1, 126 T2,
116 T3, and 190 SKIP records. A separate run-value tiering exists; the exact mapping from these
two axes to the final study shortlist needs to be frozen in the selection manifest:
\todo{reconcile the dated 83-entry shortlist, 30-task inventory, and current study suite;
record each stage's predicate, benchmark IDs, and per-benchmark inclusion/exclusion reason}.
At task and study review, unresolved reference identity, unavailable grading, incompatible
metrics or instance sets, and study-scope decisions require distinct dispositions. They must
not be collapsed into a single ``not run'' category. The reported suite can retain an explicitly
labelled protocol variant or pending result; inclusion does not certify a completed, aligned
comparison.

Practical testability records three hard blockers: grading that requires human raters at scale,
long-horizon simulations whose wall-clock requirements cannot be supported, and unavailable
proprietary access. Paid APIs, live web access, GPU compute, and token demand are recorded as
resource requirements rather than automatically treated as evidence of unsuitability. These
rules do not eliminate practical selection effects from the final evaluated suite.

\paragraph{Evidence validation.}
Benchmark-defining papers must be distinguished from papers that merely use the benchmark.
The source registry resolves identity through curated links, citation graphs, full-text search,
and web confirmation; unresolved identities enter human review. Model-extracted labels that can
exclude a candidate require a verbatim supporting quotation, validated against the retrieved
source. An absent or invalid quotation makes the field \texttt{unknown}; unknown evidence cannot
demote a benchmark. Tier assignment is deterministic given the validated fields. This safeguard
was motivated by an earlier record set in which reproducibility evidence was null in all 6{,}058
leaderboard rows even though it affected tiering.

\paragraph{Search and extraction.}
Citation retrieval combines Semantic Scholar paper identifiers, an OpenAlex citation sweep,
and full-text benchmark-name and alias search. DOI or paper identifiers deduplicate records.
Influence scores prioritize reading without excluding papers. Alias filtering precedes model
extraction; rejected passages can be sampled to evaluate recall. Low-confidence results and
candidate SOTA claims receive a stronger-model review. Full source evidence is retained for
judgment. An earlier 12{,}000-character cap truncated 96\% of the reviewed paper bodies
(median length 77{,}688 characters) and yielded 1{,}314 unknown labels; this is a pipeline
failure observation, not an estimate of current finder accuracy.

\paragraph{Score provenance and attribution.}
Every leaderboard row records method, backbone, frame, metric, split, evaluated instance count,
setting, date, quotation, and source location. Mechanical quotation checks do not guarantee
correct attribution: a historical audit found 109 of 1{,}406 rows assigned to the wrong method
after table row or column drift, although their quotations passed the mechanical check.
Candidate SOTA claims therefore require source-table attribution and an adversarial search for
a stronger claim. Official leaderboards are checked for denominator and protocol consistency.
The records distinguish \texttt{full\_set\_sota} from \texttt{best\_any\_frame}; agent,
bare-model, and classical-system reference scores are separate. Multi-agent methods are defined
operationally as two or more agent roles or instances exchanging messages, rather than a lone
tool-using agent or self-consistency procedure.

\paragraph{Finder validation still required.}
A human-adjudicated sample should evaluate benchmark identity, eligible-system retrieval,
score and method attribution, and metric/frame correctness separately. Agreement between two
model agents is not a substitute for that reference. \todo{Freeze the validation sample,
adjudication procedure, sample sizes, accuracy estimates, and uncertainty intervals.}

\section{Direct-Codex configuration and execution protocol}
\label{app:codex-setup}
\label{sec:appendix:setup}

Table~\ref{tab:codex-config} separates the settings recorded in the existing draft from settings
that need run-level provenance before the full submission. The direct-Codex arm denotes this
specific configuration, rather than every possible coding-agent configuration.

\begin{table}[ht]
\centering
\small
\begin{tabular}{@{}p{0.25\textwidth}p{0.69\textwidth}@{}}
\toprule
Field & Configuration or required record \\
\midrule
Agent & Codex CLI, version recorded per run; the original FinanceMath, ELAIPBench, and IndustryOR artifacts record 0.149.0; the ELAIPBench recovery and newer reruns record 0.149.1 \citep{openai2026codex} \\
Model endpoint & \texttt{azure/gpt-5.6-sol}; \todo{archive the provider deployment revision} \\
Reasoning effort & \texttt{medium} for the documented direct-Codex runs \\
Attempts & $k=1$ (pass@1) by default; benchmark-specific departures must be explicit \\
Scaffold & No intended domain-specific scaffold or per-benchmark tuning; instruction deviations are audited \\
Task and grader & Benchmark-specific task and declared evaluator, pinned by the run contract \\
Baseline environment & Benchmark-provided/permitted resources, each named in the prompt; superseded restricted runs are separate variants \\
Agent visibility & Instance workspace and private credentials home; gold answers and sibling results excluded \\
Network & Model-endpoint allowlist for closed-book/fixed-corpus tasks; public egress for declared web tasks \\
Provider search & Audited independently of container egress and task wording \\
Sampling & \todo{record resolved temperature, top-$p$, seed support, and endpoint defaults} \\
Resource limits & Per-benchmark turn and compute budgets in the contract; \todo{export actual token, time, and retry limits} \\
Runtime versions & \todo{export harness commit, image digests, adapter commits, and task-manifest hashes} \\
Prompts and tools & \todo{archive complete system/task prompts, emitted instructions, and tool configuration per benchmark} \\
Run provenance & \todo{freeze run IDs, execution dates, instance lists, and validity/supersession flags} \\
\bottomrule
\end{tabular}
\caption{Detailed direct-Codex setup. A missing resolved setting remains an explicit TODO;
endpoint defaults are not assumed equal to a comparator's configuration.}
\label{tab:codex-config}
\end{table}

\subsection{Container execution and checks}
\label{app:execution}

A runner adapter supplies installation and execution entry points. Each trial runs inside a
Docker container with the instance workspace and an ephemeral private credentials home. The
agent phase excludes the answer key, dataset adapter, and other trials' outputs; grading is
performed separately with access to the answer key. Comparator code must execute inside the
same isolation boundary. Running it on the host and copying only its answer into the container
does not enforce the task's environment or network restrictions.

The stage runner executes the following checks in order and stops on failure:
\begin{enumerate}
\item \texttt{manifest}: compare emitted tasks with the digests in \texttt{dataset.toml}.
\item \texttt{check}: validate task structure and whether instructions specify the required output.
\item \texttt{delivery}: confirm that staged input data reach the container.
\item \texttt{oracle}: submit a known-correct answer and verify the expected full score.
\item \texttt{nop}: check whether taking no action receives unintended credit.
\item \texttt{preflight}: verify model access through the configured network path.
\item \texttt{run}: execute the selected agent adapter.
\item \texttt{usable}: distinguish agent attempts and valid grades from sandbox or evaluator failure.
\item \texttt{analyze}: inspect whether behavior follows the task and whether instructions were sufficient.
\item \texttt{aggregate}: apply the benchmark-native headline metric and inspect the recorded traces.
\end{enumerate}

Oracle and no-action checks run after task changes without a model call. They do not replace
checking evaluator errors. Failure case C103 recorded both controls passing an evaluator whose
\texttt{pandas} import failed, because the check inspected the harness exit status instead of
trial errors. A grader that cannot run is not a measured incorrect answer. Similarly, a valid
execution does not establish freedom from training-data exposure; the run contract records a
\texttt{post\_cutoff} exposure flag but cannot verify memorization absence.

\paragraph{Null baseline and task fidelity.}
The generator summary records a constant-answer baseline. In the historical deployment snapshot,
29 of 30 benchmarks had a measured baseline; \texttt{critpt} had an ungraded preflight return,
which must not be interpreted as a measured zero. Recorded examples are FinanceMath 0.02
(4 of 200 answers within numeric tolerance), OOLONG 0.27 under its partial-credit scorer, and
SciTab 0.4074 on 108 measured instances. SciTab's separately recorded expectation over its
published task is 0.3734; the two quantities have different denominators and must not be
interchanged. These controls characterize scoring floors rather than agent capability.
Task-fidelity declarations cover search space, environment, outputs, and scope. For example,
FinanceMath's output contract was expanded from a numeric answer to the program and answer
needed to report the official accuracy and execution-rate metrics.

\paragraph{Instance and score contracts.}
Each task records both the intended evaluation count and the emission count. Prefix truncation,
failed emissions, and failed trials are handled explicitly rather than absorbed into a score's
denominator. The selected metric is validated against the exact published table column; using
an upstream scorer alone is insufficient when a paper reports a different statistic, such as
item-level F1 instead of exact match. Sampling is pass@1 unless a declared benchmark protocol
requires a different $k$. Unstated sampling in a source paper remains unstated in the record.

\paragraph{Network and tool access.}
Container egress, provider-side tools, and task instructions are checked independently.
Restricting container HTTP requests does not disable search performed by the model provider;
provider search needs its own setting. Conversely, a task with open egress can discourage all
browsing if its instructions claim that no network is available. A historical inspection found
7 of 10 traces browsing on the one web task that invited browsing, versus 0 of 9 inspected
traces across the other web tasks. This observation motivates configuration validation; it is
not a controlled estimate of browsing's effect on score. Neither network restriction prevents
information stored in model parameters.

\section{Comparator adaptation and alignment declarations}
\label{app:alignment}

\paragraph{Method versus environment.}
The baseline is determined from the benchmark's own paper or release, rather than the resources
used by the first run or a comparator-derived host list. CRAG provides its mock API; Biomni's
Eval1 uses its E1 environment; corpus-based benchmarks provide their staged corpus. Where
public web search is permitted, Codex receives a plain search tool and its prompt states the
permission. The 2026-09-12 design withdraws the old comparator-derived 15-host allowlist and
labels older restrictive settings as separate variants.

Method-specific components remain available to their own comparator: learned embeddings,
indexes, caches, fitted prompts, evolved populations, or trained checkpoints. They are not
removed to make the comparator resemble Codex or copied into the baseline. Runtime resources
provided to every entrant belong in both arms. Declared method-specific tool asymmetries support
only one-directional conclusions: Codex can match or exceed the bar, but a deficit cannot be
reported as a loss. A comparator's fitted components require a record of fitting data and any
overlap with evaluation instances.

\paragraph{Optional reference-informed arm and exposure checks.}
The reference-informed arm receives the comparator paper and a scrubbed code release, including
hand-written procedures, on top of the baseline. It tests whether Codex can use that method's
description. Scrubbing removes benchmark data, saved outputs, results, caches, tabular data,
and any files carrying instances or answers; the exact removals are documented per comparator.
An optional tool-library arm can separately expose callable method tools without the planner,
but is not part of the default study. Neither optional arm is assumed complete.
All open-web arms require scanning citations, fetched URLs, and response text for recorded gold
sources or verbatim answers; detected exposure invalidates the item and its count is reported.
Provider-side search cannot be fully constrained by container egress, so an allowlist is not a
substitute for the scan. The actual scan output must accompany a reported open-web result.

\paragraph{Protocol record.}
The benchmark's \texttt{run\_policy.json} records the evaluated split, size, emitted count,
metric, exact scoring implementation, selected published reference and source, tool access,
network mode, sampling, turn budget, comparator binding, and protocol deviations. Provenance
fields identify where web access and sampling were verified and whether the emitted metric
matches the cited published column. The comparator binding includes the system name,
repository URL, and pinned commit.

The comparator's \texttt{alignment.json} names a reference arm and declares ten variables:
instance set, answer key, grader, backbone, reasoning effort, sampling, execution environment,
network policy, tool surface, and compute budget. Each variable carries both values, evidence,
and a status of identical, equivalent, differs, or not applicable. An omitted variable is
undeclared. A deliberate method difference is marked as such; other differences require a
reported limitation or correction. \texttt{comparators/check\_alignment.py} checks the schema,
missing axes, undeclared deviations, and selected cross-field constraints. This checker cannot
establish that a declaration accurately describes executable code.

\paragraph{Code audit.}
A separate audit reads contracts against adapter code, upstream code, patches, and rendered
run artifacts. A second agent attempts to refute each finding. The 2026-09-11 audit covered
25 bindings: 15 left unassessed in the preceding synthesis and 10 without full direct-Codex
runs. It deliberately excluded five bindings previously marked verified:
\texttt{stark}, \texttt{bright}, \texttt{scitab}, \texttt{llm\_srbench}, and \texttt{industryor}.
The population and findings are reported in Section~\ref{sec:audit}. This selected audit
cannot estimate the error rate of the entire candidate corpus or deployed suite.

\paragraph{Backbone substitution and patches.}
In the historical snapshot, 28 comparator repositories were cloned and 20 had declared patch
sets; 13 patch sets were wholly or partly motivated by an endpoint that rejected an explicit
temperature. Replacing a requested temperature with endpoint-default sampling changes a method's
operating conditions, especially when it votes or repeats attempts. Patches therefore record
the upstream behavior, the change, its reason, and any supported directional implication.
The score effect is unknown without measurement. Comparator calls pass through a logging proxy,
and per-instance pricing precedes larger runs. The role-level model configuration must also be
recorded: later decisions permit lower-cost auxiliary models for some extraction, rewriting, or
page-summary roles, while benchmark judges can retain their specified models. For example, the
2026-09-15 ruling retains Test of Time's Luna extraction and Sol question answering. Thus a
shared principal backbone does not imply that every helper and judge call uses that backbone.

\paragraph{Historical reasoning and denominator caveats.}
The recorded ReLoop patch explicitly set \texttt{medium} effort. MoCA-Agent sent no effort;
the 2026-09-15 decision treats that endpoint default as medium and retains the old run, while
requiring future adapters to send and log the value explicitly. An assumed documented default
and a logged request setting are distinct evidence. The earlier AgentSPEX run recorded
\texttt{high} in all 403 workflows; a later medium-effort run was completed and its results
belong in the current comparison. The original direct-Codex snapshot completed 399 of403 trials after four pre-agent failures.
The HF recovery run re-executed those four IDs, allowing the updated paired analysis to use403.
An answer-channel asymmetry was also recorded. Separately, the historical LongBench v2
instruction supplied a hand-written retrieval strategy despite the intended absence of a domain
scaffold. Qualifications apply to the affected run IDs, not automatically to later re-runs.

\paragraph{Fixed evaluation subsets.}
The ongoing comparator campaign uses versioned manifests in
\texttt{4\_comparators/samples/}, with seed 20260917, explicit instance IDs, and recorded
sampling order. The fraction denotes approximately one quarter of each benchmark's instance
population, not one quarter of comparator systems. Examples are Biomni Eval1 (108/433),
EarthBench (61/247), and FrontierScience Research (15/60). Where declared, samples preserve
benchmark-specific strata: task family, modality subset, subject, label, or dataset subset.
Rounding and the random-generator convention are stored in each manifest; not every manifest
uses the same generator sequence. A separate draw order supports reuse of early probes in the
formal sample. Completed direct-Codex results must be restricted to these same IDs for paired
analysis; a 25\% comparator mean cannot be paired with the entire direct-Codex population.
Partial completions within the selected sample remain provisional, and final aggregation follows
the benchmark's native weighting. \todo{Freeze final coverage, completed paired IDs, missingness,
resolved settings, and confidence intervals when the ongoing runs finish.}

Two declared exceptions require separate denominators: ChemBench uses a 280-question
subset, and Test of Time reuses the mentor's existing 280-question set (five per each of 56
question-type/graph-generator cells, 10\% of 2{,}800). The Test of Time 700-question manifest
was superseded on 2026-09-17; only the ordering of its fixed 280 IDs is shuffled, using seed
20260918. Prefix pilots therefore extend to at most these 280 paired questions.

\paragraph{Interpretation of reproduction accuracy.}
A faithful reconstruction on the original configuration and an adapted evaluation on the shared
backbone are separate quantities. A difference between the adapted comparator and its published
score can reflect backbone, data, tool, sampling, or implementation changes. It cannot by itself
measure reproduction error. Where original-setting reproduction is available, report the exact
frame and residual score difference. Otherwise, report the adaptation and its unresolved
variables. \todo{Add completed 25\%-subset runs with their run IDs, common-instance coverage,
configuration checks, and separate original-setting reproduction results where available.}

\section{Review artifacts, decisions, and trace analysis}
\label{app:review}

The review interface at \texttt{2\_audit/webtool/} exposes benchmark and method records,
source evidence, and comments. Comments are anchored by benchmark identity, method identity,
and source paper, rather than table position or numeric score. This preserves review context
when records are inserted or corrected. A separate shortlist records human judgments. The
historical interface does not authenticate commenters, so comments are a review channel rather
than proof of reviewer identity.

\paragraph{Decision provenance.}
Unresolved choices belong next to the settings they govern, with a dated ruling once resolved.
The historical contracts included owner-decision fields for 12 deployed benchmarks:
\texttt{biomni\_eval1}, \texttt{bright}, \texttt{chembench}, \texttt{earthbench},
\texttt{insightbench}, \texttt{llm\_srbench}, \texttt{matscibench}, \texttt{medagentsbench},
\texttt{multihiertt}, \texttt{oolong}, \texttt{scitab}, and \texttt{test\_of\_time}.
For example, SciTab's record describes comparator optimization using labels overlapping 80 of
1{,}224 scored instances and the choices between declaring that overlap, excluding those
examples from fitting, or retaining only the paper score as an external reference. These
historical open fields must be reconciled with later rulings before finalizing the study.

\paragraph{Logs and ledgers.}
\texttt{comparators/COMPARABILITY\_CASES.md} records admissibility decisions and reversals
supported by new evidence. Its 2026-09-11 Case 16 permits published scores on other backbones
as external references, provided the distinction is labeled; it does not establish alignment
between two locally executed arms. \texttt{3\_deploy\_run/RERUNS.md} records invalid or
superseded runs, missing trials, prerequisites, and estimated recovery costs. The historical
eight-row ledger contained four invalid/superseded runs, one additional arm, and three partial
recoveries. Its old queue is not a current progress report.
\texttt{2\_audit/PIPELINE\_FAILURE\_CASES.md} records failures, observed costs, and preventive
changes. Deployment notes document the additional work needed to execute each specialist.
A correction should update the relevant pipeline stage and version stamp so that affected
records recompute, rather than alter only the displayed record.

\paragraph{Trace-analysis protocol to complete.}
The main analysis separates (i) domain-level outcome patterns, (ii) cases favoring the
specialist, and (iii) cases favoring Codex. A paired case requires the same task instance,
valid grades, and both execution traces. Candidate explanatory dimensions include access to
relevant evidence, task decomposition or representation, domain constraints, verification,
and error recovery. These dimensions are analysis questions, not established findings.
Infrastructure failures, answer-channel defects, and instruction mismatches are labeled
separately as threats to measurement. Each proposed explanation should cite the trace events
that support it and the observations that could contradict it; causal claims require an
intervention or ablation. \todo{Freeze the paired-trace sample, coding rubric, adjudication
procedure, case counts, and targeted ablations after the comparator runs complete.}
